\section{Konflik Shift-Reduce dan Reduce-Reduce}
\label{sec:ch12-konflik}

\subsection{Konflik Shift-Reduce}

Konflik shift-reduce terjadi ketika satu sel tabel memiliki aksi
shift dan reduce. Contoh klasik adalah dangling else:
\[
\text{stmt} \to \text{if}\ E\ \text{then}\ \text{stmt}
\mid \text{if}\ E\ \text{then}\ \text{stmt}\ \text{else}\
\text{stmt}.
\]
Setelah parser melihat \code{if E then stmt} dan token berikutnya
\code{else}, ia dapat menggeser \code{else} atau mengurangi
\code{stmt} terlebih dahulu.

\subsection{Konflik Reduce-Reduce}

Konflik reduce-reduce terjadi ketika dua aturan dapat direduksi pada
sel yang sama. Contohnya, jika dua nonterminal memiliki ruas kanan
yang sama dan FOLLOW yang tumpang tindih:
\[
A \to x, \qquad B \to x.
\]
Parser tidak dapat menentukan apakah \(x\) harus direduksi menjadi
\(A\) atau \(B\).

\subsection{Resolusi}

Konflik dapat diselesaikan dengan:
\begin{enumerate}[label=\arabic*.,leftmargin=*]
  \item mengubah grammar agar tidak ambigu;
  \item menambahkan precedence dan associativity;
  \item memilih aksi default secara eksplisit;
  \item memisahkan aturan yang tumpang tindih.
\end{enumerate}

Bison menyediakan deklarasi precedence:
\begin{lstlisting}[language=Bison]
%left OR
%left AND
%left EQ NE
%left LT LE GT GE
%left PLUS MINUS
%left STAR SLASH PERCENT
%right NOT UMINUS
\end{lstlisting}
Deklarasi ini memberitahu Bison cara menyelesaikan konflik
shift-reduce berdasarkan prioritas operator.

\begin{figure}[H]
\centering
\begin{tikzpicture}[node distance=6mm and 10mm]
  \node[phase, text width=3.2cm] (state) {State \(I\)};
  \node[decision, text width=2.8cm, below left=of state] (shift)
        {Shift?};
  \node[decision, text width=2.8cm, below right=of state] (reduce)
        {Reduce?};
  \node[phase, text width=3.2cm, below=of shift] (s) {Geser token};
  \node[phase, text width=3.2cm, below=of reduce] (r) {Reduksi aturan};
  \draw[arrow] (state) -- (shift);
  \draw[arrow] (state) -- (reduce);
  \draw[arrow] (shift) -- (s);
  \draw[arrow] (reduce) -- (r);
\end{tikzpicture}
\caption{Konflik shift-reduce: dua aksi untuk satu sel tabel}
\label{fig:ch12-konflik}
\end{figure}

\subsection{Membaca Laporan Konflik}

Parser generator melaporkan konflik beserta state dan token pemicu.
Laporan ini harus dibaca, bukan diabaikan. Konflik yang tidak
disengaja sering menandakan grammar yang terlalu ambigu atau aturan
yang tumpang tindih.

\begin{warningbox}
Jangan hanya menaikkan kapasitas parser atau memilih aksi default
untuk menyembunyikan konflik. Pahami penyebabnya terlebih dahulu.
Konflik yang tersembunyi dapat menghasilkan pohon urai yang salah.
\end{warningbox}
